\section{Learning with an agent}
\subsection{Understanding and cognitive debt}

\begin{frame}{Abundant output consumes scarce attention}
  \begin{center}
  \begin{tikzpicture}[
    node distance=0.52cm,
    stage/.style={draw=palegray,rounded corners=2pt,text width=2.4cm,
      minimum height=1.15cm,align=center},
    arr/.style={-{Latex[length=2mm]},thick,color=crimson}
  ]
    \node[stage] (gen) {Generate\\candidates};
    \node[stage,right=of gen] (sel) {Select\\what matters};
    \node[stage,right=of sel] (ver) {Verify\\the result};
    \node[stage,right=of ver] (exp) {Explain and\\use it};
    \draw[arr] (gen) -- (sel);
    \draw[arr] (sel) -- (ver);
    \draw[arr] (ver) -- (exp);
  \end{tikzpicture}
  \end{center}

  \vspace{0.6em}
  Faster production does not remove scarcity; it moves scarcity to the
  attention required to choose, check, explain, and connect results to the
  question.

  \vfill
  \fnote{Simon argued that information abundance creates a corresponding
  scarcity of attention \citep{simon1971attention}.}
\end{frame}

\begin{frame}{Understanding is the new bottleneck}
  \begin{columns}[T,onlytextwidth]
    \column{0.47\textwidth}
      \begin{block}{Understand to verify}
        \begin{itemize}
          \item Does the code run?
          \item Does it match the request?
          \item Do the checks pass?
          \item Should I keep it?
        \end{itemize}
      \end{block}
    \column{0.49\textwidth}
      \begin{block}{Understand to participate}
        \begin{itemize}
          \item Why was this approach chosen?
          \item Which assumptions are inside it?
          \item What question should come next?
          \item How should the project change?
        \end{itemize}
      \end{block}
  \end{columns}

  \vfill
  \fnote{Adapted and paraphrased from \citet{litt2026}. Litt argues that understanding supports approval of the current work and active participation in the next project loop.}
\end{frame}

\begin{frame}{Cognitive debt}
  The project can continue while the people responsible for it lose the ability
  to explain:
  \begin{itemize}
    \item what the system is supposed to do;
    \item why an important decision was made;
    \item which assumptions the output carries;
    \item how to change the work safely.
  \end{itemize}

  \vspace{0.45em}
  This is \alert{cognitive debt}: the work advances faster than shared
  understanding. The cost appears later, when we need to modify, teach, or
  repair the project.

  \vfill
  \fnote{Concept adapted from the discussion in \citet{litt2026}, which connects the term to Margaret-Anne Storey and Simon Willison.}
\end{frame}

\begin{frame}{External assistance is not practical judgment}
  \begin{columns}[T,onlytextwidth]
    \column{0.47\textwidth}
      \begin{block}{Plato: reminder is not understanding}
        In the \emph{Phaedrus}, writing is presented as an external aid that can
        preserve words while giving a learner only the appearance of wisdom.

        \vspace{0.45em}
        \alert{Student test:} Can I question, explain, and reconstruct what the
        agent produced?
      \end{block}
    \column{0.49\textwidth}
      \begin{block}{Aristotle: technique is not judgment}
        \emph{Techn\={e}} concerns making something well. \emph{Phron\={e}sis},
        or practical wisdom, concerns deliberating well about situations that
        could be otherwise.

        \vspace{0.45em}
        \alert{Policy test:} What should be done, for whom, and under which
        uncertainty?
      \end{block}
  \end{columns}

  \vfill
  AI can extend memory and technique. Education still has to cultivate
  understanding and practical judgment.

  \fnote{Paraphrased from Plato, \emph{Phaedrus} 274c--275b, and Aristotle,
  \emph{Nicomachean Ethics}, Book VI \citep{platoPhaedrus,aristotleEthics}.}
\end{frame}

\begin{frame}{Checkpoint 1: Is this understanding?}
  \alert{Block 2 · minutes 70--74}

  \vspace{0.55em}
  An agent produces a clean plot. The student can describe its appearance but
  cannot explain which rows were excluded.

  \vspace{0.55em}
  Decide together:
  \begin{enumerate}
    \item What evidence of cognitive debt do you see?
    \item What must the student do before keeping the plot?
  \end{enumerate}

  \vfill
  \begin{center}
    \textbf{2 minutes discuss} \hspace{2.0em} $\longrightarrow$ \hspace{2.0em}
    \textbf{2 minutes share}
  \end{center}
\end{frame}

\begin{frame}{Education produces an artifact and a future capability}
  \begin{align*}
    q_t &= q(h_t,a_t),\\
    h_{t+1} &= h_t + \text{practice}
      + \text{AI scaffolding} - \text{displaced practice}.
  \end{align*}

  \begin{itemize}
    \item $q_t$: the quality of today's script, table, figure, or explanation.
    \item $h_t$: the student's current ability to understand and change it.
    \item $a_t$: the form and intensity of AI assistance.
  \end{itemize}

  \vfill
  Assistance can improve today's artifact while either building or weakening
  tomorrow's capability. The learning design determines which effect dominates.
\end{frame}

\begin{frame}{What one coding experiment suggests}
  \citet{shen2026} randomly assigned AI assistance to 52 developers learning a
  new Python library.

  \begin{columns}[T,onlytextwidth]
    \column{0.48\textwidth}
      \begin{block}{Immediate product}
        The AI group finished about two minutes earlier. The difference was not
        statistically significant.
      \end{block}
    \column{0.48\textwidth}
      \begin{block}{Later mastery}
        The AI group scored 50\% on the quiz, compared with 67\% without AI. The
        largest gap appeared in debugging.
      \end{block}
  \end{columns}

  \vfill
  This small study examines one skill. We use its result as a warning and avoid
  generalizing it to every learning setting.
\end{frame}

\begin{frame}{AI can transmit patterns without transmitting judgment}
  \begin{columns}[T,onlytextwidth]
    \column{0.47\textwidth}
      \begin{block}{Polanyi: knowledge can be tacit}
        Practice contains patterns, exceptions, and judgment that experts may
        use without being able to state as a complete rule
        \citep{polanyi1966tacit}.
      \end{block}
    \column{0.49\textwidth}
      \begin{block}{Workplace evidence}
        AI assistance raised customer-support productivity by 14\% on average
        and by 34\% among novice and lower-skilled workers
        \citep{brynjolfsson2023work}.
      \end{block}
  \end{columns}

  \vspace{0.55em}
  One interpretation is that the system helped distribute patterns previously
  concentrated among experienced workers. The open question is whether users
  also learn when those patterns do not apply.

  \vfill
  Coding-agent traces still show persistent returns to task expertise; that
  evidence is observational rather than causal \citep{hitzig2026}.
\end{frame}

\subsection{A learning routine}

\begin{frame}{Before, during, and after agent help}
  \small
  \begin{tabularx}{\textwidth}{@{}>{\bfseries}p{2.0cm}p{4.6cm}X@{}}
    \toprule
    Moment & Student action & Useful agent support \\
    \midrule
    Before & Predict the result and state the criterion. & Ask for a plan or a small example. \\
    During & Compare the proposal with your prediction. & Ask why and test a small case. \\
    After & Reconstruct and explain the central decision. & Ask for a quiz or an independent check. \\
    \bottomrule
  \end{tabularx}

  \vfill
  The goal is to keep our understanding close to the speed of production
  \citep{shen2026,litt2026}.
\end{frame}

\begin{frame}{The student standard}
  Before keeping a substantial block of agent-written code, be ready to state:
  \begin{enumerate}
    \item its purpose;
    \item its input and output;
    \item one important assumption;
    \item one consequence of changing it;
    \item one check that could reveal a mistake.
  \end{enumerate}

  \vfill
  If this is difficult, ask the agent for a smaller example, an explanation with
  background, or a short quiz. Then try again in your own words.
\end{frame}

\begin{frame}{Use the agent to help you understand}
  \begin{columns}[T,onlytextwidth]
    \column{0.31\textwidth}
      \begin{block}{Explainer}
        Ask for the goal, background, important choices, and a literate walkthrough.
      \end{block}
    \column{0.31\textwidth}
      \begin{block}{Quiz}
        Ask five questions that test whether you can reconstruct the central change.
      \end{block}
    \column{0.31\textwidth}
      \begin{block}{Micro-example}
        Ask for a small object you can change and observe. Another long explanation may be harder to test.
      \end{block}
  \end{columns}

  \vfill
  These suggestions adapt Litt's techniques to a classroom: explanation,
  retrieval practice, and small environments that the learner can manipulate
  \citep{litt2026}.
\end{frame}

\begin{frame}{Checkpoint 2: Design a learning loop}
  \alert{Block 2 · minutes 84--88}

  \vspace{0.55em}
  An agent proposes code that joins the indicator dictionary to the WDI data.

  \vspace{0.55em}
  Choose one action for each moment:
  \begin{enumerate}
    \item What will you predict \emph{before} running it?
    \item What will you inspect \emph{during} the run?
    \item What will you explain or test \emph{afterward}?
  \end{enumerate}

  \vfill
  \begin{center}
    \textbf{2 minutes discuss} \hspace{2.0em} $\longrightarrow$ \hspace{2.0em}
    \textbf{2 minutes share}
  \end{center}
\end{frame}
